\paragraph{Separating approximation error from numerical precision.}
To determine whether the observed discrepancies arise from the quadrature
approximation or its numerical evaluation, we compare GPU FP32, CPU FP32,
and CPU FP64 explanations for the same fitted Cox models and empirical
background games. We retain the 30 fixed training-background patients,
five distinct held-out patients per dataset, and automatic node choices for
$\varepsilon\in\{10^{-1},10^{-2},10^{-3}\}$.
The GPU results are reused from the saved experiments; both CPU controls
are newly measured. Inputs, numerical sources, and package versions were
verified against the GPU experiment. The CPU controls evaluate the same
parallel product-game formula at the same selected node counts. CPU FP32
provides a control on the same hardware as CPU FP64, so that changing the
device is not necessary to investigate the effect of precision.

Across the three tolerances, mean CPU FP64 latencies are 156--164\,ms
for GSE24080 and 992--1,035\,ms for TCGA LGG, compared with
85--90\,ms and 310--320\,ms for the saved GPU FP32 runs.
These are measured precision/device tradeoffs, not predictions of
FP64 GPU performance.

All three configurations are compared with the same independent 64-node
FP64 reference $\phi^{\mathrm{ref}}_p$:
\[
\Delta_a=\max_{p=1,\ldots,5}
\|\widehat\phi^{a}_{m_p,p}-\phi^{\mathrm{ref}}_p\|_\infty,
\qquad a\in\{\mathrm{GPU32},\mathrm{CPU32},\mathrm{CPU64}\}.
\]
This reference has analytical quadrature bounds below $10^{-10}$ and a
96-node cross-check for the first patient of each dataset. It is a
high-accuracy numerical reference, not a full degree-exact FP64 computation.
In particular, $\Delta_{\mathrm{GPU}}$ in
Table~\ref{tab:cox-cpu-gpu-precision} now uses this common reference;
it is not the discrepancy against rounded GPU degree-exact output reported
in the earlier table. All automatically selected rules satisfy their
analytical bounds $C_p\le\varepsilon$ in real arithmetic.

Against the common reference, GPU FP32 meets the requested threshold in
19/30 patient--tolerance cases, CPU FP32 in 19/30,
and CPU FP64 in 30/30. The worst errors over these cases are
$3.91\times10^{-2}$, $3.98\times10^{-2}$, and $1.59\times10^{-10}$, respectively.
The reduction in error from CPU FP32 to CPU FP64, with the device and
formula held fixed, supports limited work precision as an important source
of the discrepancies. Differences between the CPU FP32 and GPU FP32
outputs additionally show the role of backend-specific evaluation.
These numerical controls support this explanation but do not constitute
a formal floating-point error bound.

For timing, each distinct dataset/node-count configuration is warmed up
once before its first measured call. We record one synchronized API call
per patient and method, including automatic node selection, data movement,
and host accumulation. Model fitting, compilation/warm-up, and cached rule
construction are excluded. The CPU and GPU runs use an 8\,GiB planning
budget and one background row per block; this is a planning budget, not
measured memory consumption. Reported standard deviations describe
variation across patients rather than repeated-call uncertainty.
The saved GPU timings and new CPU timings were recorded in separate runs.
The CPU FP64/GPU FP32 latency ratio changes both device and precision,
so it should not be interpreted as a matched-precision GPU speedup.
Full degree-exact CPU runtimes were not measured and are left blank.

The analytical guarantee remains a bound on quadrature error for the
fixed fitted model and chosen 30-row empirical background game.
It excludes floating-point rounding, uncertainty in the fitted model,
and error from replacing a population background with this empirical
background. Small FP64 discrepancies are an empirical accuracy check,
not a stronger analytical certificate. The observed accuracy of CPU FP64
must therefore be paired with its CPU runtime rather than with the faster
GPU FP32 runtime.
